\documentclass[10pt,a4paper,oneside]{book}
\usepackage[a4paper,top=22mm,bottom=21mm,left=22mm,right=20mm,headheight=15pt]{geometry}
\usepackage{fontspec}
\usepackage{polyglossia}
\setdefaultlanguage{thai}
\setotherlanguage{english}
\setmainfont{LINE Seed Sans TH}
\setsansfont{LINE Seed Sans TH}
\newfontfamily\thaifont{LINE Seed Sans TH}
\setmonofont{Menlo}[Scale=0.82]
\newfontfamily\thaifonttt{LINE Seed Sans TH}
\usepackage[table]{xcolor}
\definecolor{brand}{HTML}{0052FF}
\definecolor{navy}{HTML}{0F172A}
\definecolor{muted}{HTML}{64748B}
\definecolor{pale}{HTML}{E8F4FF}
\definecolor{line}{HTML}{E2E8F0}
\usepackage{hyperref}
\hypersetup{colorlinks=true,linkcolor=brand,urlcolor=brand,pdftitle={Elegant Solution - คู่มือการใช้งาน การติดตั้ง และการบริหารจัดการระบบ},pdfauthor={Elegant Solution},pdfsubject={Complete System, Installation and Technical Manual},pdfkeywords={Elegant Solution, Tracking Agent, คู่มือ}}
\usepackage{fancyhdr,booktabs,longtable,tabularx,array,enumitem,caption,tikz,listings}
\usetikzlibrary{arrows.meta,positioning,shapes.geometric}
\setlist{nosep,leftmargin=18pt,topsep=5pt}
\setlength{\parindent}{0pt}
\setlength{\parskip}{5pt}
\renewcommand{\arraystretch}{1.22}
\setlength{\tabcolsep}{2.5pt}
\newcolumntype{L}[1]{>{\raggedright\arraybackslash}p{#1}}
\def\UrlBreaks{\do\/\do\-\do\_\do\.\do\;}
\Urlmuskip=0mu plus 1mu
\emergencystretch=3em
\tolerance=2000
\hbadness=10000
\pagestyle{fancy}
\fancyhf{}
\fancyhead[L]{\small\color{muted}Elegant Solution \quad / \quad คู่มือระบบ 0.5.0}
\fancyhead[R]{\small\color{muted}\nouppercase{\leftmark}}
\fancyfoot[L]{\small\color{muted}ฉบับตรวจจากโค้ด 12 พ.ย. 2026}
\fancyfoot[R]{\small\color{muted}\thepage}
\renewcommand{\headrulewidth}{0.3pt}
\renewcommand{\footrulewidth}{0pt}
\fancypagestyle{plain}{\fancyhf{}\fancyfoot[R]{\small\color{muted}\thepage}\renewcommand{\headrulewidth}{0pt}}
\renewcommand{\chaptername}{บทที่}
\renewcommand{\contentsname}{สารบัญ}
\renewcommand{\figurename}{ภาพที่}
\renewcommand{\tablename}{ตารางที่}
\newcommand{\code}[1]{{\small\texttt{\detokenize{#1}}}}
\newcommand{\source}[1]{\par{\footnotesize\color{muted}หลักฐานใน repository: \nolinkurl{#1}}\par}
\newcommand{\status}[1]{{\color{brand}\bfseries #1}}
\newcommand{\note}[1]{\par\noindent\colorbox{pale}{\parbox{0.955\linewidth}{\textbf{หมายเหตุ}\quad #1}}\par}
\newcommand{\warn}[1]{\par\noindent\colorbox{line}{\parbox{0.955\linewidth}{\textbf{ข้อควรระวัง}\quad #1}}\par}
\newcommand{\step}[1]{\item #1}
\lstset{basicstyle=\ttfamily\footnotesize,breaklines=true,breakatwhitespace=false,frame=single,rulecolor=\color{line},backgroundcolor=\color{pale!25},columns=fullflexible,keepspaces=true,showstringspaces=false,aboveskip=7pt,belowskip=7pt}
% Editable original vector diagrams. Keep labels short enough for A4.
\newcommand{\flowfour}[5]{%
\begin{figure}[htbp]\centering
\begin{tikzpicture}[node distance=3mm,box/.style={draw=line,fill=pale,rounded corners=4pt,align=center,text width=31mm,minimum height=17mm,font=\small},arr/.style={-{Latex[length=2mm]},draw=brand,line width=1pt}]
\node[box] (a) {#2}; \node[box,right=of a] (b) {#3}; \node[box,right=of b] (c) {#4}; \node[box,right=of c] (d) {#5};
\draw[arr] (a) -- (b); \draw[arr] (b) -- (c); \draw[arr] (c) -- (d);
\end{tikzpicture}\caption{#1}\end{figure}}
\newcommand{\flowthree}[4]{%
\begin{figure}[htbp]\centering
\begin{tikzpicture}[node distance=6mm,box/.style={draw=line,fill=pale,rounded corners=4pt,align=center,text width=42mm,minimum height=17mm,font=\small},arr/.style={-{Latex[length=2mm]},draw=brand,line width=1pt}]
\node[box] (a) {#2}; \node[box,right=of a] (b) {#3}; \node[box,right=of b] (c) {#4};
\draw[arr] (a) -- (b); \draw[arr] (b) -- (c);
\end{tikzpicture}\caption{#1}\end{figure}}

\begin{document}
\frontmatter
\begin{titlepage}
\pagecolor{pale}\color{navy}
\vspace*{25mm}
\begin{center}
\begin{tikzpicture}[scale=0.82]
  \fill[brand!22,rounded corners=17pt] (0,0) rectangle (3.2,4.1);
  \draw[brand,line width=7pt,line cap=round] (0.7,1.5) .. controls (1.5,0.6) and (2.5,1.7) .. (3.6,2.8);
  \draw[brand!80!black,line width=4pt,line cap=round] (1.05,3) -- (1.05,3.3);
  \draw[brand!80!black,line width=4pt,line cap=round] (2.25,3) -- (2.25,3.3);
\end{tikzpicture}
\vspace{15mm}
{\fontsize{31}{38}\selectfont\bfseries Elegant Solution\par}
\vspace{9mm}
{\fontsize{18}{27}\selectfont คู่มือการใช้งาน การติดตั้ง\par และการบริหารจัดการระบบ\par}
\vspace{7mm}
{\large Complete System, Installation \& Technical Manual\par}
\vfill
{\color{brand}\rule{0.68\textwidth}{1.5pt}\par}
\vspace{6mm}
{\large Version 0.5.0 \quad | \quad 12 พฤศจิกายน 2026\par}
{\small อ้างอิง repository ที่ commit \code{ada86dadbaa1e4711cca99c56c1f1d6d83b2174a}\par}
\end{center}
\end{titlepage}
\nopagecolor
\chapter*{วิธีอ่านคู่มือนี้}
\addcontentsline{toc}{chapter}{วิธีอ่านคู่มือนี้}
คู่มือนี้ใช้โค้ดใน branch \code{feature/cross-platform-tracking-agent-exam} ณ commit ที่ระบุบนปกเป็นหลักฐานหลัก เอกสารประวัติและผลทดสอบที่ผู้ดูแลรายงานแยกจากพฤติกรรมที่ยืนยันในโค้ด ทุกตัวอย่างใช้ข้อมูลสมมติ คำสั่งที่แตะ production ต้องผ่านขั้นอนุมัติและสำรองข้อมูลตามบทที่ 9

\textbf{ป้ายสถานะ}: \status{IMPLEMENTED} = มีโค้ด; \status{AUTOMATED TESTED} = มีผลตรวจอัตโนมัติที่รายงาน; \status{VERIFIED ON REAL HARDWARE} = มีผลจากเครื่องจริงที่ผู้ใช้รายงาน; \status{NOT YET VERIFIED} = ยังไม่มีหลักฐาน; \status{PARTIAL IMPLEMENTATION} = มีโค้ดบางส่วน; \status{NOT IMPLEMENTED} = ไม่พบการทำงานนั้นในโค้ดที่ตรวจ

\note{ผลทดสอบ macOS ล่าสุดมาจากคำบอกและผล Terminal ที่ผู้ใช้ส่งในบทสนทนา: ติดตั้ง/activate, reboot แล้ว UI เปิดเอง, เครื่องเดิม Online ไม่มีรายการซ้ำ, จากนั้นถอนด้วย maintenance code และตรวจว่า receipt/service/UI หายไป การ login นักศึกษาและ foreground activity ยังไม่มีผลทดสอบบนเครื่องจริงในหลักฐานชุดนี้}
\tableofcontents
\mainmatter
\part{คู่มือผู้ใช้}
\chapter{รู้จัก Elegant Solution}
Elegant Solution เป็นระบบจัดการการใช้ห้องปฏิบัติการที่เชื่อมบัญชีผู้ใช้ ตารางเรียน การค้นหาห้อง การจอง สิทธิ์เข้าห้อง จุดตรวจหน้าห้อง และสถานะเครื่องคอมพิวเตอร์ ผู้ใช้ทำงานผ่านเว็บ; ผู้ดูแลกำหนดห้อง อนุมัติคำขอ ตรวจประวัติ และติดตั้ง Tracking Agent บนเครื่องแล็บ จุดมุ่งหมายของโครงการคือให้ข้อมูลการใช้ห้องและเครื่องสัมพันธ์กันในระบบเดียว โดยไม่อ้างผลสำรวจหรือการทดลองที่ไม่มีหลักฐาน

ขอบเขตสำคัญมีสองชั้น: เว็บและ Face AI ทำงานบนเซิร์ฟเวอร์; Tracking Agent ทำงานบนเครื่องแล็บแต่ละเครื่อง การมี Agent Online หมายถึงเครื่องส่ง heartbeat ได้ ไม่ได้แปลว่ามีนักศึกษาผ่านการตรวจสิทธิ์แล้ว และซอฟต์แวร์ยืนยันการเข้าห้องไม่ได้พิสูจน์ว่ากลอนไฟฟ้าปลดล็อกจริง
\flowfour{ภาพรวมระบบ (แผนภาพอธิบายจากโค้ด)}{บัญชีและตารางเรียน}{ค้นหา/จองห้อง}{Gate และ Face AI}{Agent และรายงาน}
\section*{ความสามารถหลักและขอบเขต}
ระบบมีเว็บสำหรับ \code{STUDENT}, \code{GUEST}, \code{ADMIN}, \code{SUPER_ADMIN}; ตารางเรียน, OCR, Room Finder, booking, token QR, face enrollment/verification, Admin Tracking และ batch enrollment. หน้า Contact/Survey เป็นการแสดงผลบางส่วนและไม่ควรอ้างว่ารับข้อมูลจริงโดยไม่ตรวจโค้ดปลายทาง เว็บไซต์ blacklist ใน Agent รุ่น 0.5.0 ไม่ได้ยืนยันการบังคับใช้ครบทุก OS. การบังคับ login ก่อนใช้งาน desktop macOS ยังไม่ใช่พฤติกรรมของ Agent ปัจจุบัน
\source{web/src/lib/permissions.ts; web/src/lib/room-entitlement.ts; agent/src/session_ui.py}

\chapter{บทบาทและสิทธิ์}
สิทธิ์ไม่ได้มาจากชื่อหน้าที่ UI แสดงอย่างเดียว Backend ตรวจ session, role, account status และในกรณีห้อง ตรวจ entitlement จากข้อมูลจริง การ login เว็บ, การยืนยันตัวเครื่องด้วย API key และ Agent user session เป็นคนละสิ่ง
\begin{longtable}{L{34mm}L{31mm}L{31mm}L{31mm}L{31mm}}
\toprule รายการ & SUPER\_ADMIN & ADMIN & STUDENT & GUEST \\ \midrule\endhead
เว็บและโปรไฟล์ & ได้ & ได้ & ได้ & ได้เมื่อ login \\
เมนู Admin ยกเว้น System Control & ได้ & ได้ & ไม่ได้ & ไม่ได้ \\
System Control และ Deployment Batch & ได้ & ไม่ได้ & ไม่ได้ & ไม่ได้ \\
Room Finder & ดูได้ & ดูได้ & ดูได้ & เฉพาะ approved และห้องที่เปิดให้ Guest \\
จองห้องด้วยตนเอง & ไม่ได้ & ไม่ได้ & ได้ & เฉพาะ approved \\
ลงทะเบียนใบหน้า & ได้ & ได้ & ได้ & ไม่ได้ \\
Agent user login & มีบัญชี approved และสิทธิ์ห้อง ณ เวลาใช้งาน & เช่นเดียวกัน & เช่นเดียวกัน & เช่นเดียวกัน แต่ guest ไม่มี face enrollment \\
\bottomrule
\end{longtable}
\flowfour{ชั้นการตรวจสิทธิ์ (แผนภาพอธิบาย)}{Web session}{Role/สถานะบัญชี}{สิทธิ์ห้องตามเวลา}{Machine + Agent session}
การสมัครผ่านหน้าเว็บสร้าง \code{GUEST/PENDING}; เจ้าหน้าที่อนุมัติเป็น \code{APPROVED} หรือปฏิเสธเป็น \code{REJECTED}. นักศึกษา Google ใหม่ใช้โดเมน \code{bumail.net} ที่ยืนยันอีเมลแล้ว และถูกสร้างเป็น \code{STUDENT/APPROVED}. บัญชี ADMIN/SUPER\_ADMIN ที่มีอยู่ใช้ Google ได้เมื่ออีเมลยืนยัน โดย role ยึดฐานข้อมูล. ดูสิทธิ์หน้า Admin จาก server action/API ไม่ใช่รายการเมนูอย่างเดียว
\source{web/src/lib/permissions.ts; web/src/lib/auth.ts; web/src/lib/google-sign-in.ts; web/src/lib/admin-nav.ts; web/src/app/api/agent/auth/route.ts}

\chapter{สมัครบัญชีและเข้าสู่ระบบ}
\section*{สมัคร Guest}
\begin{enumerate}
\step{เปิด \code{/register} กรอกชื่อ อีเมล รหัสผ่าน และเบอร์โทรถ้ามี; หน้าฟอร์มตรวจรหัสผ่านอย่างน้อย 6 ตัวอักษรและยืนยันซ้ำ}
\step{ส่งแบบฟอร์ม ระบบตรวจอีเมลซ้ำและเก็บรหัสผ่านแบบ bcrypt; บัญชีใหม่มี role \code{GUEST} สถานะ \code{PENDING}}
\step{รอผู้ดูแลตรวจที่ \code{/admin/users}. หากยังไม่อนุมัติ Guest ใช้ Room Finder/booking ไม่ได้; หากถูกปฏิเสธให้ติดต่อผู้ดูแล}
\end{enumerate}
\section*{เข้าสู่ระบบและจัดการบัญชี}
ที่ \code{/login} ใช้อีเมล/รหัสผ่าน หรือ Google ตามกฎโดเมนข้างต้น หากข้อมูลผิดจะไม่สร้าง web session. หน้า \code{/profile} แก้ชื่อ/โทรศัพท์และเปลี่ยนรหัสผ่านด้วยรหัสเดิมได้ นักศึกษาผูก student ID 10 หลักครั้งแรก; หากต้องแก้ภายหลังให้ติดต่อ admin. ไม่พบ workflow ส่งอีเมล reset password อัตโนมัติในโค้ดที่ตรวจ การ logout เว็บไม่ใช่การ logout Agent session บนเครื่องแล็บ
\warn{บัญชีที่ login ด้วย Google และมีรหัสผ่านสุ่มภายใน ไม่ควรอ้างว่าผู้ใช้รู้รหัสผ่านนั้น การเปลี่ยนรหัสผ่านที่หน้าโปรไฟล์ต้องผ่านรหัสปัจจุบันตามโค้ด}
\source{web/src/actions/auth.ts; web/src/lib/auth.ts; web/src/actions/profile.ts; web/src/app/(auth)/register/page.tsx}

\chapter{คู่มือนักศึกษา}
\section*{Dashboard และตารางเรียน}
หลัง login เข้า \code{/dashboard} เพื่อดูข้อมูลห้อง ตารางเรียน และรายการที่เกี่ยวข้องกับบัญชี ตารางเรียนสัมพันธ์กับภาคการศึกษาที่ระบบกำหนด; ถ้าไม่พบรายการให้ตรวจ semester และ student ID ก่อน ระบบใช้วัน/เวลาแบบ Asia/Bangkok ในการให้สิทธิ์ห้อง ผู้ใช้ควรตรวจความถูกต้องของวิชา ห้อง วัน และช่วงเวลาเสมอ
\section*{นำเข้าตาราง BU Mobile ด้วย OCR}
\begin{enumerate}
\step{เตรียมภาพหน้าจอตารางเรียนที่ตัวอักษรชัด ครบชื่อวิชา วัน เวลา และห้อง; เปิดแผง OCR ใน Dashboard}
\step{เลือกภาพ ระบบโหลด \code{tesseract.js@5.1.1} ใน browser แล้วอ่านข้อความจากภาพ}
\step{ตรวจ preview ที่แยกรายการพร้อมคำเตือน/ข้อผิดพลาด; ตรวจ course code, section, ชื่อวิชา, วัน, slot และห้อง}
\step{แก้ข้อมูลที่จำเป็นใน preview ตาม UI แล้วสั่ง import เฉพาะรายการที่พร้อม ระบบรวมรายการซ้ำที่ parser พบ; ตรวจตารางที่บันทึกแล้วอีกครั้ง}
\end{enumerate}
\note{OCR เป็นการช่วยอ่านภาพ ไม่ใช่หลักฐานว่าข้อมูลถูกต้องทุกช่อง ข้อความดิบและ parser debug มีไว้ตรวจปัญหาใน UI; อย่าใช้ภาพตารางจริงของผู้อื่นในคู่มือหรือการสาธิต}
\section*{ค้นหาห้องและจอง}
เปิด Room Finder จาก Dashboard เลือกวันที่, slot, จำนวนผู้ใช้, ความสามารถที่ต้องการ, อาคารและชั้นตามตัวกรองที่ UI มี ผล \code{Available Now} ขึ้นกับช่วงเวลาปัจจุบัน ตารางเรียน การจอง หน้าต่างเปิดให้จอง การปิดห้อง และความจุ ห้องที่เปิด booking แต่ยังไม่มี availability window จะถูกปฏิเสธแบบ fail-closed

เลือกห้องและ slot $\rightarrow$ ตรวจผู้เข้าร่วม $\rightarrow$ ยืนยันคำขอ. การจองเดี่ยวคือผู้เข้าร่วมหนึ่งคน; booking แบบกลุ่มใช้ snapshot รายชื่อสมาชิก ณ เวลาสร้าง booking. สถานะ \code{PENDING} ต้องรออนุมัติ; \code{APPROVED} ให้สิทธิ์ตามช่วงเวลาที่อนุมัติ การสร้าง booking ยังไม่สร้าง session, access token หรือสั่งประตูเปิด หน้าสร้าง Group QR ใน Dashboard ถูกปิดด้วย \code{ENABLE_GROUP_UI = false} แม้มีโค้ดกลุ่มอยู่
\flowfour{การจองห้องของนักศึกษา (แผนภาพอธิบาย)}{ค้นหาวันและ slot}{เลือกห้อง/ผู้เข้าร่วม}{ตรวจเงื่อนไข}{PENDING หรือ APPROVED}
\section*{เข้าห้องและประวัติ}
สิทธิ์ปัจจุบันมาจาก \code{CLASS}, booking ที่อนุมัติและมีชื่อผู้ใช้ใน participant snapshot, หรือ \code{LEGACY_TOKEN} ที่ออกให้ตนเอง. CLASS เริ่มก่อนเวลาเรียน 30 นาที. ที่ Gate สามารถสแกน QR token หรือใช้ face verification ตามช่องทางที่เปิดใช้; successful face check-in/check-out บันทึก Session. QR token เดิมมี validation ของตนเองตามบท 20; อย่าสมมติว่า QR ทุกอันให้สิทธิ์ผ่าน entitlement engine แบบเดียวกัน. ตรวจประวัติใน Dashboard/รายงานตามสิทธิ์
\source{web/src/components/ocr-timetable-panel.tsx; web/src/lib/availability.ts; web/src/lib/room-booking.ts; web/src/lib/room-entitlement.ts; web/src/app/(app)/dashboard/dashboard-client.tsx}

\chapter{คู่มือ Guest}
Guest สมัครที่ \code{/register} แล้วรออนุมัติจาก Admin. หลังสถานะ \code{APPROVED} จึงค้นหาและจองห้องที่ตั้ง \code{guestAccess=true} และไม่ใช่ \code{staffOnly}. การเลือกเวลาและอนุมัติ booking ใช้หลักเดียวกับนักศึกษา แต่ Guest ไม่มีเส้นทางลงทะเบียนใบหน้าใน \code{/api/face/enroll}; การสาธิต Gate ด้วยใบหน้าของ Guest จึงไม่ใช่ workflow ที่รองรับเต็มรูปแบบในรุ่นนี้

\begin{enumerate}
\step{ตรวจสถานะบัญชีและ login เว็บ}
\step{เปิด Room Finder; ถ้าห้องไม่ปรากฏ ให้ Admin ตรวจ \code{guestAccess}, \code{staffOnly}, เวลาเปิด และความจุ}
\step{ส่ง booking และรออนุมัติหากนโยบายห้องเป็น MANUAL}
\step{ดูสถานะคำขอและติดต่อ Admin ถ้าถูกปฏิเสธ; อย่าใช้ QR หรือบัญชีของผู้อื่นแทน}
\end{enumerate}
\source{web/src/lib/permissions.ts; web/src/actions/auth.ts; web/src/app/api/face/enroll/route.ts; web/src/lib/room-booking.ts}

\chapter{คู่มือ Admin}
เมนูของ Admin ประกอบด้วย ภาพรวม, จัดการผู้ใช้, จัดการห้อง, คำขอจองห้อง, ตารางจองห้อง, นำเข้าตารางเรียน, ติดตามการใช้งาน, System Health, ประวัติการใช้งาน, รายงาน และวิเคราะห์ข้อมูล. หน้า \code{/admin/system} สงวนให้ SUPER\_ADMIN ตามการนำทางและ server action

\section*{งานประจำวัน}
\begin{enumerate}
\step{เปิด \code{/admin/dashboard} ตรวจภาพรวม แล้วไป \code{/admin/users} เพื่ออนุมัติ Guest/จัดการบัญชีตามสิทธิ์}
\step{ที่ \code{/admin/rooms} ตรวจรหัสห้อง สถานะ ความจุ capabilities, guest access, staff-only, นโยบาย \code{AUTO/MANUAL}, availability windows และ closures}
\step{ที่ \code{/admin/bookings} ตรวจคำขอ \code{PENDING}, ผู้ร่วมจอง, วัน/slot และความขัดแย้งก่อนอนุมัติหรือปฏิเสธ; ดู \code{/admin/calendar} ประกอบ}
\step{ที่ \code{/admin/schedule} นำเข้าตารางเรียนตามภาคเรียน และตรวจแถวผิดพลาดก่อนใช้งาน Gate}
\step{ที่ \code{/admin/tracking} ดูเครื่อง Online, ผู้ใช้ปัจจุบัน และประวัติกิจกรรม; Online ของเครื่องไม่ยืนยันว่ามีนักศึกษา login Agent}
\step{ตรวจ \code{/admin/diagnostics}, \code{/admin/logs}, \code{/admin/reports}, \code{/admin/analytics} เพื่อวิเคราะห์เหตุการณ์และส่งออกข้อมูลตามที่หน้าเว็บรองรับ}
\end{enumerate}
Admin Tracking มีเมนูสร้าง maintenance code ผูกเครื่องและ action; code มีอายุ 5 นาทีและใช้ครั้งเดียว การแก้ website blacklist ผ่านหน้า Tracking เป็นการตั้งค่า backend; การบังคับใช้จริงขึ้นกับ Agent/OS และยังไม่ควรอ้างว่าครบทุกแพลตฟอร์ม
\source{web/src/lib/admin-nav.ts; web/src/actions/agent.ts; web/src/app/admin/tracking/tracking-client.tsx; web/src/lib/room-admin.ts}

\chapter{คู่มือ Super Admin}
SUPER\_ADMIN มีเมนู Admin ทั้งหมดและหน้า \code{/admin/system} เพิ่มเติม ใช้ตั้งค่าระบบและจัดการ Tracking Agent Deployment. การสร้าง batch ต้องเลือกห้องที่ active และจำนวนเต็ม 1--100 เครื่องต่อ batch. CSV ที่ดาวน์โหลดได้ครั้งเดียวควรเก็บในพื้นที่จำกัดสิทธิ์และแจกหนึ่ง code ต่อหนึ่งเครื่อง
การตั้งค่า SUPER\_ADMIN คนแรกเป็นขั้นตอนที่ต้องจัดการโดยผู้ดูแลระบบที่ได้รับมอบหมาย: seed script ที่ตรวจสร้างบัญชีตัวอย่าง role \code{ADMIN} ไม่ได้สร้าง SUPER\_ADMIN. การสร้าง/เลื่อนสิทธิ์ผู้ดูแลผ่านเว็บต้องมี SUPER\_ADMIN อยู่ก่อน จึงไม่ควรอ้างว่ามี bootstrap อัตโนมัติหรือรัน seed ตัวอย่างบน production
\begin{enumerate}
\step{เปิด \code{/admin/system} ส่วน Tracking Agent Deployment; เลือกห้องและจำนวนเครื่อง}
\step{สร้าง batch; ตรวจ room code, จำนวนที่ร้องขอ และเวลา expiry}
\step{ดาวน์โหลด CSV ใน browser ทันที เก็บอย่างปลอดภัย; แจก code ตาม machine code ที่จองไว้}
\step{ติดตั้ง PKG/EXE/DEB ให้ตรง OS; เปิด UI และกรอก code}
\step{กลับมาดูสถานะ \code{PENDING}, \code{ACTIVATED}, \code{REVOKED}, \code{EXPIRED}; ตรวจ Admin Tracking เพิ่มเติม}
\end{enumerate}
Batch code มีอายุ 7 วัน; code เดี่ยวจาก Admin Tracking มีอายุ 10 นาที. การ revoke batch ปิดเฉพาะ code ที่ยังไม่ใช้ ไม่ได้เพิกถอน machine ที่ activate ไปแล้ว. Installer download ต้องตั้ง \code{AGENT_INSTALLER_DIR} และ stage artifact บน backend ก่อน หน้า System Control ยังมีตัวเลือกการตั้งค่าระบบอื่นตามโค้ด; เปลี่ยนค่าเฉพาะเมื่อทราบผลกระทบและบันทึกเหตุผล
\source{web/src/actions/agent-deployments.ts; web/src/lib/agent-deployment-batches.ts; web/src/lib/agent-activation.ts; web/src/app/api/admin/agent/installers/[platform]/route.ts; web/prisma/seed.ts}

\part{การติดตั้งและบริหารระบบ}
\chapter{สถาปัตยกรรมติดตั้งใช้งาน}
Production ที่เคยตรวจบนเซิร์ฟเวอร์ Ubuntu 20.04.6 ใช้ Cloudflare เป็น HTTPS สาธารณะ, Nginx proxy ไป Next.js พอร์ต 3000, PM2 ชื่อ \code{elegant}, PostgreSQL 16 ใน Docker container \code{elegant-postgres} และ Face AI แยกพอร์ต 8000. นี่เป็น topology ที่บันทึกไว้ใน runbook; คู่มือนี้ไม่ได้ SSH เข้าเซิร์ฟเวอร์เพื่อยืนยันสถานะสด
\flowfour{เส้นทาง production ที่บันทึกไว้ (แผนภาพสถาปัตยกรรม)}{Browser/Agent\ HTTPS}{Cloudflare\ Nginx}{Next.js\ PM2}{PostgreSQL\ Face AI}
เครื่องแล็บ Windows/macOS/Ubuntu รัน Tracking Agent และติดต่อ \code{/api/agent/*} ทาง HTTPS. Backend เก็บแพ็กเกจ installer ใน directory ภายนอก web root เพื่อให้ Admin ดาวน์โหลด; \textbf{ไม่ต้องติดตั้ง Tracking Agent บนเซิร์ฟเวอร์ backend}. Machine service กับ UI อยู่บนเครื่องปลายทาง
\section*{เครือข่ายและขอบเขต}
ตรวจ DNS/TLS ไป \url{https://elegant.fafromh0me.com}; Nginx ส่งต่อ \code{localhost:3000}; PostgreSQL ผูก localhost. Runbook ระบุ Face AI พอร์ต 8000 บน server และการตรวจเดิมพบ origin พอร์ต 3000/8000 เข้าถึงจากภายนอกได้ จึงต้องตรวจ firewall/reverse proxy และวางแผนจำกัดการเข้าถึงใน change แยก Agent ใช้เฉพาะ HTTPS backend ไม่ต้องติดต่อ PostgreSQL หรือ Face AI โดยตรง
\source{docs/PRODUCTION_PM2_DEPLOYMENT.md; agent/README.md; web/src/app/api/agent/heartbeat/route.ts}

\chapter{Backend: ติดตั้งและ deploy}
\section*{ก่อน deploy}
สภาพ production ที่บันทึกล่าสุดมี Node.js \code{v24.13.1}, npm \code{11.8.0}, PM2 \code{5.3.0}, global pnpm \code{8.8.0}; release ใหม่ใช้ \code{npx --yes pnpm@9.15.9} แบบแยก. RAM ว่างประมาณ 2.4 GiB, swap ถูกใช้ 1.8/4 GiB และ disk ว่าง 11 GB ณ preflight; ค่าจริงต้องตรวจซ้ำ. Repository production มีไฟล์แก้และ untracked จึงใช้ Git worktree/release directory ใหม่ ไม่ overwrite checkout เดิม

\begin{lstlisting}[language=bash]
# ตรวจอ่านอย่างเดียวบน server ก่อนเริ่มขั้นเปลี่ยนระบบ
node --version
pm2 list
free -h
df -h /home/lamy
sudo docker exec elegant-postgres sh -lc \
  'pg_isready -U "$POSTGRES_USER" -d "$POSTGRES_DB"'
git -C /home/lamy/Elegant-Solutions status --short
\end{lstlisting}
ตั้ง release จาก full commit SHA ที่อนุมัติแล้วตาม \code{docs/PRODUCTION_PM2_DEPLOYMENT.md}. คัดลอก runtime \code{.env} แบบรักษาสิทธิ์ 600 และตรวจด้วย \code{cmp -s}; ห้ามพิมพ์ค่า secret. ตรวจ \code{DATABASE_URL}, \code{AUTH_SECRET}, \code{AUTH_URL} และ environment ที่เกี่ยวข้องว่ามีค่าและไม่ขัดกับ PM2 เดิมโดยไม่แสดงค่า. \code{AGENT_INSTALLER_DIR} ชี้ directory ของ artifacts ซึ่งไม่อยู่ใน web root
\warn{ห้ามรัน \code{web/prisma/seed.ts} กับ production เพื่อสร้าง Super Admin: script นี้สร้างบัญชีตัวอย่างด้วย credential ที่ฝังไว้และไม่ได้สร้าง SUPER\_ADMIN. การเตรียมบัญชีผู้ดูแลคนแรกต้องทำตามกระบวนการอนุมัติและข้อมูลจริงของหน่วยงาน}

\section*{ลำดับหลังอนุมัติ production change}
\begin{enumerate}
\step{สำรอง PostgreSQL ด้วย \code{pg_dump -Fc}; ตรวจ checksum และ \code{pg_restore --list}; ถ้าเป็นไปได้ restore ทดสอบสำเนาแยก}
\step{สร้าง isolated release จาก commit ที่ตรวจ migration history แล้ว; ติดตั้ง dependencies ด้วย pnpm 9.15.9; รัน Prisma generate และ build โดยจำกัด worker/heap ตาม runbook}
\step{ตรวจ \code{prisma migrate status}. หากแสดง pending เฉพาะรายการที่คาดจึงรัน \code{prisma migrate deploy}; ห้ามใช้ reset/db push แบบทำลายข้อมูล}
\step{เริ่ม build ใหม่แบบ standalone บนพอร์ต 3001 เพื่อตรวจ \code{/api/liveness}; หลังผ่านจึงสลับ PM2 ตาม runbook และตรวจ HTTPS public endpoint}
\step{บันทึก SHA, PM2 cwd/command, สถานะ migrations และ health checks; เก็บ release เดิมเพื่อ rollback}
\end{enumerate}
\flowfour{การปล่อย backend (แผนภาพอธิบาย)}{สำรอง DB}{Build release แยก}{Migrate + smoke 3001}{สลับ PM2/ตรวจผล}
\warn{บล็อกต่อไปเป็นตัวอย่างคำสั่งสำหรับ release ที่แยกแล้ว ไม่ใช่คำสั่งให้รันทันทีบน production. ค่าภายใน \code{.env} และ backup ห้ามใส่ใน Git/คู่มือ}
\begin{lstlisting}[language=bash]
cd "$RELEASE/web"
NODE_ENV=development npx --yes pnpm@9.15.9 install --frozen-lockfile
npx --yes pnpm@9.15.9 exec prisma generate --config ./prisma.config.ts
CIRCLE_NODE_TOTAL=2 NODE_OPTIONS=--max-old-space-size=1536 \
  npx --yes pnpm@9.15.9 build
# หลังอนุมัติและมี backup ที่ตรวจแล้วเท่านั้น:
npx --yes pnpm@9.15.9 exec prisma migrate status --config ./prisma.config.ts
npx --yes pnpm@9.15.9 exec prisma migrate deploy --config ./prisma.config.ts
\end{lstlisting}
คำสั่งสำรองและตรวจ archive บน server ใช้ \code{sudo docker exec} ตามสิทธิ์ที่ตรวจแล้ว. การอ่านรายการใน archive ไม่เท่ากับการ restore ทดสอบ; ก่อน cutover ต้องมีผล restore แยกในสภาพแวดล้อมที่ควบคุมได้
\begin{lstlisting}[language=bash]
# PRODUCTION CHANGE: ทำหลังอนุมัติเท่านั้น
BACKUP_DIR=/home/lamy/elegant-db-backups
install -d -m 700 "$BACKUP_DIR"
BACKUP="$BACKUP_DIR/elegant_$(date +%Y%m%d_%H%M%S).dump"
umask 077
sudo docker exec elegant-postgres sh -lc \
  'pg_dump -U "$POSTGRES_USER" -d "$POSTGRES_DB" -Fc' > "$BACKUP"
pg_restore --list "$BACKUP" >/dev/null
sha256sum "$BACKUP"
\end{lstlisting}
การ start build ใหม่ใช้ Node รัน standalone \code{server.js} ไม่พึ่ง global pnpm 8. Smoke test ที่พอร์ต 3001 ก่อน แล้วจึง cutover พอร์ต 3000 ผ่าน PM2. ตรวจจาก localhost และ public HTTPS; คำสั่งจริงต้องใช้ \code{STANDALONE_SERVER} และ \code{STANDALONE_DIR} ที่ตรวจแล้วจาก release ตาม runbook
\begin{lstlisting}[language=bash]
# PRODUCTION CHANGE: เฉพาะหลัง migrate/build/smoke ผ่าน
pm2 stop elegant
PORT=3000 HOSTNAME=0.0.0.0 pm2 start "$STANDALONE_SERVER" \
  --name elegant-v050 --interpreter "$(command -v node)" \
  --cwd "$STANDALONE_DIR"
curl --fail --silent --show-error \
  http://127.0.0.1:3000/api/liveness
curl --fail --silent --show-error \
  https://elegant.fafromh0me.com/api/liveness
pm2 list
pm2 save
\end{lstlisting}
Rollback แอปคือกลับไป PM2 release เดิมตาม runbook; schema ใหม่ถูกออกแบบให้ backward compatible กับแอปเดิม แต่การกู้ DB คือขั้นแยกที่ต้องวางแผน downtime/ข้อมูลใหม่หลัง backup. อย่าทับฐานข้อมูลสดด้วย dump อัตโนมัติ อ่านคำสั่ง PM2 cutover และ rollback แบบเต็มจาก runbook ก่อนปฏิบัติ
\begin{lstlisting}[language=bash]
# APP ROLLBACK: หลังอนุมัติแก้เหตุขัดข้อง
pm2 stop elegant-v050 >/dev/null 2>&1 || true
pm2 restart elegant
curl --fail --silent --show-error \
  https://elegant.fafromh0me.com/
pm2 save
\end{lstlisting}
หากต้องกู้ฐานข้อมูล ให้ restore backup ลงฐาน \emph{ใหม่} ก่อน เปรียบเทียบ migration/จำนวน records และขออนุมัติแยกก่อนเปลี่ยน \code{DATABASE_URL}. การ restore ทับ production โดยตรงหลังมี writes ใหม่ทำข้อมูลหลังเวลาสำรองสูญหาย
\source{docs/PRODUCTION_PM2_DEPLOYMENT.md; web/next.config.ts; web/prisma/schema.prisma}

\chapter{Tracking Agent 0.5.0}
Agent แยก \code{machine_service.py} จาก \code{session_ui.py}. Service เก็บ machine identity/API key ในพื้นที่ระบบ ส่ง heartbeat ทุก 10 วินาทีแม้ไม่มีนักศึกษา login และ sync config. UI ใน graphical session ใช้ code เปิดใช้งาน, login ผู้ใช้ และอ่าน foreground app/title เฉพาะเมื่อ Agent user session active. UI ติดต่อ service ผ่าน loopback \code{127.0.0.1:18767} โดย local bridge token; service ส่งข้อมูลไป backend
\flowfour{องค์ประกอบ Agent (แผนภาพอธิบาย)}{UI ใน desktop}{Loopback + token}{Machine service}{HTTPS API}
\section*{วงจร session และความเป็นส่วนตัว}
Activation ใช้ code หกหลักครั้งเดียวและได้ machine code/API key ซึ่งบันทึกแบบสิทธิ์จำกัด. Agent login ตรวจ email/password บัญชี \code{APPROVED} และ entitlement ของห้องที่ผูกเครื่อง; backend ออก session token ที่ผูก user, machine, room และเวลาสิ้นสุด. เมื่อ session หมดอายุ service refresh สิทธิ์หรือ clear session. ไม่ส่งชื่อหน้าต่างโดยไม่มี session. Backend รับ machine-only heartbeat แล้วอัปเดต \code{lastSeen}; \code{AgentLog} สร้างเมื่อมี session ที่ใช้ได้
\flowfour{Heartbeat และกิจกรรม (แผนภาพลำดับ)}{Boot service}{Machine heartbeat}{Agent login}{Activity + heartbeat}
\note{Agent login ปัจจุบันไม่ล็อก macOS/Windows/Linux desktop ก่อนใช้งาน เครื่องยังใช้ได้โดยไม่ login Agent. นี่เป็นข้อจำกัดของการควบคุมการใช้เครื่อง และช่วงนั้นไม่มี activity ที่ผูกผู้ใช้}
\source{agent/src/machine_service.py; agent/src/session_ui.py; agent/src/config.py; web/src/app/api/agent/heartbeat/route.ts; web/src/app/api/agent/auth/route.ts}

\chapter{Batch Enrollment และรหัสเปิดใช้งาน}
SUPER\_ADMIN สร้าง Deployment Batch ที่ \code{/admin/system}. Backend ตรวจห้อง active และจำนวน 1--100, จอง machine code รูปแบบ \code{<roomCode>-PC<sequence>}, สร้าง code หกหลักด้วย \code{randomInt}, hash ก่อนเก็บ และกำหนดหมดอายุ 7 วัน. Transaction ลองแก้ collision ตามจำนวนครั้งที่กำหนด. CSV ที่ browser ได้ครั้งเดียวมีคอลัมน์ \code{batch_id}, \code{room_code}, \code{machine_code}, \code{activation_code}, \code{expires_at} ตามลำดับ
\flowfour{Batch Enrollment (แผนภาพอธิบาย)}{เลือกห้อง/จำนวน}{สร้าง batch + codes}{แจก CSV ครั้งเดียว}{activate รายเครื่อง}
\begin{enumerate}
\step{เลือกห้องและจำนวนเครื่องใน System Control; ตรวจเลขที่จองและวันหมดอายุ}
\step{เก็บ CSV ในที่จำกัดสิทธิ์; แจกหนึ่งรหัสต่อหนึ่งเครื่อง ไม่ส่งผ่านช่องทางสาธารณะ}
\step{ติดตั้งแพ็กเกจ OS ที่ถูกต้องแล้วกรอก code ใน UI; ตรวจ \code{machineCode/roomCode}}
\step{ดู deployment progress ใน System Control และสถานะ Online ใน Admin Tracking}
\end{enumerate}
รหัสถูกใช้ได้ครั้งเดียว. การ revoke batch/code มีผลกับ code ที่ยังไม่ถูกใช้; เครื่องที่ activate แล้วต้อง deactivate/uninstall แยก. Failed attempts ถูกบันทึกและมี durable activation throttle. หาก reinstall หลังถอนซึ่งลบ config ต้องออก code ใหม่; ไม่ใช้ code เก่าซ้ำ
\source{web/src/lib/agent-deployment-batches.ts; web/src/lib/agent-activation.ts; web/src/lib/agent-activation-throttle.ts; web/src/app/api/agent/activate/route.ts; web/src/app/admin/system/system-client.tsx}

\chapter{ติดตั้งและถอน Agent บน Windows}
เป้าหมาย Windows 10/11 x64. Installer คือ \code{ElegantAgentSetup_0.5.0_x64.exe} จาก Admin download หรือ artifact ที่ตรวจ SHA-256 แล้ว. Build บน Windows ใช้ Python 3.12, PyInstaller และ Inno Setup 6; ไม่ใช่ MSI. Installer ลง service \code{ElegantSolutionsMachine} และ UI auto-start ใน HKLM Run เมื่อผู้ใช้ Windows login
\begin{enumerate}
\step{รับ EXE จากผู้ดูแล; เทียบ hash กับ release ที่แจกอย่างปลอดภัย; login ด้วยบัญชี Administrator}
\step{Run as administrator; รอ service แล้วเปิด UI. กรอก activation code ของเครื่องนี้}
\step{ตรวจ service/Run registration ด้วย PowerShell ด้านล่าง; ตรวจเครื่อง Online ใน Admin Tracking ภายในประมาณ 10--30 วินาที}
\step{ทดสอบ Agent login ด้วยบัญชีที่มี entitlement ปัจจุบัน, foreground activity, logout, machine-only heartbeat และ reboot}
\end{enumerate}
\begin{lstlisting}
Get-Service ElegantSolutionsMachine
sc.exe qc ElegantSolutionsMachine
Get-ItemProperty 'HKLM:\Software\Microsoft\Windows\CurrentVersion\Run' -Name ElegantAgentUI
Test-Path "$env:ProgramData\ElegantSolution\Agent\config.json"
\end{lstlisting}
Program files อยู่ \code{C:\Program Files\ElegantSolution\Agent}; config/API key/logs อยู่ \code{C:\ProgramData\ElegantSolution\Agent}. ถ้า UI ไม่ขึ้น ให้ตรวจ registry Run และ Windows Event Viewer. \textbf{ยังไม่มีหลักฐานติดตั้ง Windows จริง} จากชุดข้อมูลนี้
\section*{ถอนและอัปเดต}
ใน Admin Tracking สร้าง maintenance code ชนิด \code{UNINSTALL_AGENT} ของเครื่องนั้นก่อน แล้วเปิด \code{C:\Program Files\ElegantSolution\Agent\uninstall.bat} แบบ Administrator. Script ตรวจ code กับ backend, หยุด/ลบ service และ registry, ลบ process/งานเก่าที่ระบุ, ลบ hosts entries เฉพาะส่วนที่ Agent จัดการ แล้วลบ ProgramData/Program Files ของ Agent. Backend deactivate machine เมื่ออนุญาต code; หากขั้นลบในเครื่องล้มเหลวหลังยืนยัน code ให้หยุดและตรวจสถานะ ไม่ใช้คำสั่งลบสุ่ม. อัปเดตโดย installer ใหม่ทับ; config ใน ProgramData ไม่อยู่ใน payload
\source{agent/build.bat; agent/setup.iss; agent/setup.bat; agent/uninstall.bat; agent/WINDOWS_INSTALLATION.md}

\chapter{ติดตั้งและถอน Agent บน macOS}
เลือก PKG ให้ตรง \code{uname -m}: \code{arm64} สำหรับ Apple Silicon หรือ \code{x86_64} สำหรับ Intel. PKG รุ่นพัฒนา \code{0.5.0} ยัง unsigned/unnotarized; ให้ผู้ดูแลประเมินแหล่งที่มาและ hash ก่อนอนุญาตเฉพาะแพ็กเกจที่เชื่อถือ อย่าปิด Gatekeeper ทั้งระบบ. การแจกจ่ายถาวรควรลงลายเซ็นและ notarize
\begin{lstlisting}[language=bash]
uname -m
shasum -a 256 -c elegant-agent_0.5.0_darwin_$(uname -m).pkg.sha256
sudo installer -pkg elegant-agent_0.5.0_darwin_$(uname -m).pkg -target /
sudo launchctl print system/com.elegantsolution.agent.machine
launchctl print gui/$(id -u)/com.elegantsolution.agent.session
\end{lstlisting}
PKG วาง binary ที่ \code{/Library/Application Support/ElegantSolution/bin/}; LaunchDaemon ทำงานตอน boot, LaunchAgent เริ่ม UI เมื่อ login desktop; config/API key ใน \code{/Library/Application Support/ElegantSolution/Agent/}. หลัง activate ตรวจ Admin Tracking $\rightarrow$ Agent login ที่มี entitlement $\rightarrow$ activity $\rightarrow$ logout $\rightarrow$ reboot. การอ่านชื่อ app/window ผ่าน System Events ต้องได้สิทธิ์ Automation/Accessibility; ถ้าไม่ได้สิทธิ์ activity อาจเป็น null แต่ heartbeat ต่อได้
\section*{หลักฐานเครื่องจริงและถอน}
ผู้ใช้รายงานว่า macOS PKG activate สำเร็จ, หลัง reboot UI เปิดเองไม่ต้อง activate ซ้ำ, Admin Tracking เห็นเครื่องเดิม Online โดยไม่ login นักศึกษาและไม่มีรายการซ้ำ. ต่อมารัน uninstaller แล้วได้ \code{Elegant Agent removed}; \code{pkgutil --pkg-info} ไม่พบ receipt และ \code{launchctl print} ไม่พบสอง service. นี่คือหลักฐานจากผู้ใช้ ไม่ใช่การตรวจเครื่องโดยผู้เขียนคู่มือโดยตรง. ผล Agent login/activity ยังไม่ยืนยัน
\begin{lstlisting}[language=bash]
# สร้าง UNINSTALL_AGENT maintenance code ใน Admin Tracking ก่อน
sudo '/Library/Application Support/ElegantSolution/bin/uninstall-agent.sh'
pkgutil --pkg-info com.elegantsolution.agent
sudo launchctl print system/com.elegantsolution.agent.machine
launchctl print gui/$(id -u)/com.elegantsolution.agent.session
\end{lstlisting}
Uninstaller ตรวจ code แล้วหยุด LaunchDaemon/LaunchAgent ลบ plist, binary, \code{Agent} และ \code{Bridge}, ลบ PKG receipt. ไม่ลบไฟล์ \code{/var/log/elegant-agent-service.log} และ \code{/var/log/elegant-agent-service.err} ตาม script. Backend deactivate machine เมื่อรับ code ถอนสำเร็จ; หาก local script ล้มเหลวภายหลัง อาจเกิดสถานะ deactivate ก่อนลบครบ. Reinstall หลังถอนต้องใช้ activation code ใหม่; update PKG ทับเก็บ config
\source{agent/build-macos.sh; agent/macos/pkg-scripts/postinstall; agent/macos/uninstall-agent.sh; agent/MACOS_INSTALLATION.md; web/src/app/api/agent/maintenance/authorize/route.ts}

\chapter{ติดตั้งและถอน Agent บน Ubuntu}
เป้าหมาย Ubuntu 22.04/24.04 amd64. DEB เป็น frozen binary; machine service ใช้ \code{elegant-agent.service} ผ่าน systemd และ UI ใช้ XDG autostart หลัง graphical login. X11 อาจอ่าน foreground ผ่าน EWMH; Wayland ส่ง app/title เป็น null เมื่อ API ไม่รองรับ จึงต้องทดสอบ desktop จริง
\begin{lstlisting}[language=bash]
uname -m
sha256sum -c elegant-agent_0.5.0_amd64.deb.sha256
sudo apt install ./elegant-agent_0.5.0_amd64.deb
systemctl status elegant-agent.service --no-pager
systemctl is-enabled elegant-agent.service
journalctl -u elegant-agent.service -n 50 --no-pager
\end{lstlisting}
เข้า desktop แล้วกรอก activation code ใน UI. ตรวจ Online, Agent login/activity, logout/heartbeat และ reboot. Config/API key อยู่ \code{/var/lib/elegant-agent/config.json}; token เชื่อม UI อยู่ \code{/run/elegant-agent/token}; binary ที่ \code{/usr/lib/elegant-agent/}. \textbf{ยังไม่มีหลักฐานติดตั้ง Ubuntu จริง} ในข้อมูลที่ตรวจ
\section*{ถอน อัปเดต และ legacy}
สร้าง \code{UNINSTALL_AGENT} maintenance code ใน Admin Tracking แล้วรัน \code{sudo elegant-agent-uninstall}; script ตรวจ code แล้ว \code{apt-get purge -y elegant-agent}. ตรวจ \code{systemctl status elegant-agent.service}, \code{dpkg -s elegant-agent} และไฟล์ config ตามนโยบาย package postrm. อัปเดตด้วย \code{sudo apt install ./<new-version>.deb}. การย้ายจาก user-level agent เดิมมี \code{sudo elegant-agent-migrate-legacy <desktop-username>} ซึ่งคัดลอก config เฉพาะเมื่อปลายทางว่างและปิด autostart เก่า; สำรอง config อย่างปลอดภัยก่อน
\source{agent/build-linux.sh; agent/linux/elegant-agent.service; agent/linux/elegant-agent.desktop; agent/linux/uninstall-agent.sh; agent/linux/postrm; agent/LINUX_INSTALLATION.md}

\chapter{บำรุงรักษาและแก้ปัญหา}
ใช้รูปแบบ อาการ $\rightarrow$ สาเหตุที่เป็นไปได้ $\rightarrow$ ตรวจ $\rightarrow$ แก้ $\rightarrow$ ยืนยัน. เริ่มจากการตรวจแบบอ่านอย่างเดียวก่อนเสมอ
\begin{longtable}{L{31mm}L{40mm}L{42mm}L{43mm}}
\toprule อาการ & สาเหตุ & ตรวจ & แก้และยืนยัน \\ \midrule\endhead
UI ไม่เปิด & graphical login/autostart ไม่ทำงาน & launchctl/Run/XDG entry & เปิด UI ตามคู่มือ OS แล้วทดสอบ login ใหม่ \\
Machine service หยุด & service startup/permission & \code{launchctl}, \code{Get-Service}, \code{systemctl} & ตรวจ log, restart ด้วย maintenance code, ดู heartbeat \\
Activate ไม่ได้ & code ผิด/หมดอายุ/ใช้แล้ว, URL/TLS & code status, HTTPS, response 403/429 & ออก code ใหม่ให้เครื่องถูกต้อง; รอตาม Retry-After \\
เครื่อง Offline & heartbeat หยุด, backend/เน็ตล่ม & service status, DNS/TLS, \code{/api/liveness} & กู้การเชื่อมต่อ; รอ heartbeat รอบใหม่ \\
Agent login ถูกปฏิเสธ & บัญชีไม่ approved/ไม่มี entitlement & ตรวจ role, room, CLASS/BOOKING/token และเวลา & แก้ข้อมูลต้นเหตุแล้ว login ใหม่ \\
ไม่มีชื่อแอป & ไม่ login Agent, macOS permission, Wayland & ตรวจ session/UI/permission & ให้สิทธิ์ OS ตามนโยบาย; ตรวจ activity ใหม่ \\
Maintenance code ถูกปฏิเสธ & ผิด action/เครื่อง, หมด 5 นาที, ใช้แล้ว & ตรวจรายการใน Admin Tracking & สร้าง code ใหม่โดยไม่ deactivate เครื่องก่อน \\
Checksum ไม่ตรง & artifact ผิด/ดาวน์โหลดเสีย & เทียบ SHA-256 จาก release ที่เชื่อถือ & หยุดติดตั้ง ดาวน์โหลดใหม่ \\
\bottomrule
\end{longtable}
\warn{หาก uninstall ได้รับอนุญาตจาก backend แล้วแต่ลบไฟล์ในเครื่องไม่สำเร็จ ให้บันทึก error และสถานะ service ก่อนดำเนินการต่อ เพราะ backend อาจ deactivate machine ไปแล้ว. อย่าเผยแพร่ config, API key, CSV หรือ log ที่มีข้อมูลบุคคล}
\source{agent/src/maintenance_cli.py; agent/src/api_client.py; web/src/app/api/agent/activate/route.ts; web/src/app/api/agent/maintenance/authorize/route.ts}

\part{อ้างอิงเทคนิคสำหรับผู้พัฒนา}
\chapter{สถาปัตยกรรมและเทคโนโลยี}
เว็บใช้ Next.js 16.1.6, React 19.2.3, NextAuth 5 beta, Prisma 7.3 และ PostgreSQL 16 ตาม package/schema/runbook. Face AI เป็น FastAPI + InsightFace รุ่นโมเดลที่กำหนดใน service. OCR ใช้ Tesseract.js 5.1.1 ที่ browser ผ่าน CDN และ parser ใน TypeScript. Tracking Agent เป็น Python core แยก service/UI; PyInstaller package สำหรับ Windows, macOS และ Linux. CI สร้างสี่ artifacts แต่ CI build ไม่ใช่หลักฐาน end-to-end บนเครื่องจริง
\flowfour{องค์ประกอบซอฟต์แวร์ (แผนภาพอธิบาย)}{Next.js UI/API}{Prisma/PostgreSQL}{Face AI}{Agent หลาย OS}
เส้นทาง request: Browser login/web actions ผ่าน NextAuth และ server action; Gate ส่ง QR/ภาพไป API; Next.js ตรวจข้อมูล/สิทธิ์ฐานข้อมูล; Face API เรียก service ภายใน. Agent ใช้ machine key แยกจาก web session. PM2/Nginx/Cloudflare อยู่ในชั้น deploy ไม่ใช่โค้ด business logic
\source{web/package.json; web/prisma/schema.prisma; web/ai-service/main.py; agent/README.md; .github/workflows/agent-packages.yml}

\chapter{ฐานข้อมูลและความสัมพันธ์}
Prisma ใช้ \code{User} เป็นบัญชี/role/status; \code{Room} เป็นห้องและนโยบาย; \code{ClassSchedule} ผูก user/room/semester; \code{RoomBooking} ผูกผู้ขอ/ห้อง/วัน/status โดย \code{RoomBookingSlot} ระบุช่วง และ \code{RoomBookingParticipant} บันทึก snapshot ผู้มีสิทธิ์. \code{AccessToken} คือ legacy token ช่วงเวลา; \code{Session} บันทึก check-in/out; \code{FaceProfile} เก็บ embedding ของ user

ด้าน Agent มี \code{AgentMachine} ผูกห้องและ machine code/API key; \code{AgentSession} ผูก machine/user/room พร้อม token hash และ expiry; \code{AgentLog} บันทึกกิจกรรม; \code{AgentActivationCode}, \code{AgentDeploymentBatch}, \code{AgentMachineSequence}, \code{AgentActivationThrottle}, \code{AgentMaintenanceCode} และ \code{AgentAuditEvent} จัดการวงจรเปิดใช้และตรวจสอบย้อนหลัง
\flowfour{ER แบบย่อ: ห้องและการใช้งาน (ภาพอธิบาย)}{User + Room}{ClassSchedule\ RoomBooking}{AccessToken\ Session}{AgentMachine\ AgentSession/Log}
Constraint ที่ควรรู้: \code{Room.roomCode} และ \code{AgentMachine.machineCode} เป็น unique; booking participant unique ต่อคู่ booking/user; slot unique ต่อ booking/slot; availability window unique ต่อ room/day/slot. การลบห้อง/ผู้ใช้มีข้อจำกัดจาก FK; ห้ามลบข้อมูลจริงเพื่อแก้ migration อย่างรวดเร็ว. ชุด migration ที่เคยซ้อมมี 17 รายการ และ production upgrade เคยรายงาน 17/17; คู่มือนี้ไม่ได้อ่านฐานข้อมูล production สด
\source{web/prisma/schema.prisma; docs/PRODUCTION_PM2_DEPLOYMENT.md; web/src/lib/room-entitlement.ts}

\chapter{กระบวนการ OCR ตารางเรียน}
Browser โหลด Tesseract.js จาก CDN, สร้าง worker, อ่านภาพแต่ละไฟล์เป็นข้อความ จากนั้น \code{parseOcrBatch} แยกบรรทัด/กลุ่มวิชา ตรวจชื่อ, semester, day, course code, section, class type, เวลา, room และเทียบ academic slot. Parser รวมแถวซ้ำและให้สถานะ \code{ready/warning/error}. UI แสดง preview และ raw OCR/debug เพื่อให้ผู้ใช้ตรวจ ก่อน import ผ่าน server action; server ตรวจข้อมูลอีกครั้งก่อนบันทึก
\flowfour{OCR pipeline (ภาพอธิบาย)}{Screenshot}{Tesseract.js}{Parser/validation}{Preview $\rightarrow$ DB}
ตัวอย่างสมมติ: ภาพมีข้อความ \code{CS101 LAB Mon 09:00-12:00 LAB-A}; preview ควรแสดงรหัสวิชา, ประเภท LAB, วัน/เวลา และห้องที่อ่านได้ ผู้ใช้ต้องแก้ถ้าภาพเบลอหรือ parser จับ slot ไม่ได้. ต้องทดสอบด้วยภาพที่ได้รับอนุญาต; OCR ไม่ตรวจความจริงของตารางจากมหาวิทยาลัยโดยตรง
\source{web/src/components/ocr-timetable-panel.tsx; web/src/lib/ocr-timetable.ts; web/src/actions/ocr-timetable.ts; web/src/lib/academic-slots.ts}

\chapter{Face AI และการยืนยันใบหน้า}
\section*{Enrollment}
\code{POST /api/face/enroll} ต้องมี web session และปฏิเสธ GUEST. UI เก็บภาพจากกล้อง 1--3 ภาพ; API ส่ง \code{user_id} และภาพ base64 ไป Face AI \code{/enroll}. Service ใช้ InsightFace ตรวจหน้า เลือกใบหน้าที่ detection score สูงสุดในแต่ละภาพ ตรวจขนาด/คุณภาพ และเฉลี่ย embedding ที่ได้. Next.js upsert \code{FaceProfile} และบันทึก log. โค้ดที่ตรวจไม่ได้ยืนยันว่าภาพมีหน้าเดียวเสมอ; การเลือกใบหน้าที่คะแนนสูงสุดอาจไม่ตรงคนที่ตั้งใจ
\flowfour{Face enrollment (ภาพอธิบาย)}{Web session + กล้อง}{Face AI detect}{Embedding เฉลี่ย}{FaceProfile}
\section*{Verification}
\code{POST /api/face/verify} รับภาพกับ \code{room_code}; server ตรวจห้อง active, ส่ง embeddings ของ non-guest users ให้ Face AI เปรียบเทียบ 1:N ด้วย threshold ค่าเริ่ม \code{0.73} หรือ \code{FACE_THRESHOLD}. เมื่อ match แล้วโหลด user, ตรวจ \code{staffOnly}, และสำหรับ STUDENT/GUEST ตรวจ entitlement ห้องปัจจุบัน; ADMIN/SUPER\_ADMIN ไม่ผ่าน entitlement check นี้ตามโค้ด. หากอนุญาตจะเรียก \code{checkInAction}; ถ้ามี active session เดิมจึง \code{checkOutAction}. Log เก็บผลการยืนยัน
\flowfour{Face verification (ภาพอธิบาย)}{Gate ส่งภาพ/room}{Face AI match}{Entitlement}{Check-in/out}
ภาพไม่ถูกต้อง, ไม่พบหน้า, ไม่มี profile หรือ service ล่มต้องปฏิเสธ/คืน error ตาม API. ไม่พบ liveness หรือ anti-spoofing ใน flow ที่ตรวจ; ไม่อ้างว่ากันภาพถ่ายได้. การตรวจผ่านซอฟต์แวร์ไม่ยืนยันกลอนไฟฟ้าเปิดจริง. Embedding เป็นข้อมูลชีวมิติ ต้องจำกัดสิทธิ์และกำหนดนโยบายการเก็บ/ลบ
\source{web/src/app/api/face/enroll/route.ts; web/src/app/api/face/verify/route.ts; web/ai-service/main.py; web/src/actions/sessions.ts}

\chapter{QR, access token และ Gate}
Dashboard สร้าง token ด้วย \code{crypto.randomUUID()} ผ่าน \code{generateAccessTokenAction} แล้วใช้ library \code{qrcode} สร้างภาพ QR จาก \emph{ค่า token} โดยตรง ไม่ใช่ URL. Gate page \code{/gate/[roomCode]} ใช้ \code{html5-qrcode} อ่าน QR แล้วเรียก server action \code{verifyGateAccessAction} โดยตรง; route \code{POST /api/gate/verify} เป็นทางเข้า API แยกที่เรียก action เดียวกัน. Validation ตรวจ token ตรง room และอยู่ในช่วง \code{validFrom/validTo}; ถ้าไม่ใช่ token ที่ใช้ได้จะลองตีความ input เป็น \code{studentId} เพื่อหาตารางเรียนของห้อง/เวลานั้น
\flowfour{QR/token flow (ภาพอธิบาย)}{Dashboard สร้าง UUID}{QR แสดง token}{Gate สแกน}{API ตรวจ room/เวลา}
\section*{ความต่างที่ต้องไม่สับสน}
\begin{itemize}
\item การสร้าง QR เป็นการแสดง token; ไม่เท่ากับการสร้าง booking และไม่ใช่ code ใช้ครั้งเดียวโดยอัตโนมัติ
\item \code{verifyAccessTokenAction} ตรวจช่วงเวลาและ room; token สามารถใช้ซ้ำในช่วงเวลาที่กำหนดตามโค้ดที่ตรวจ ไม่พบการ consume หลังสแกน
\item \code{verifyGateAccessAction} มีเส้นทาง student ID fallback; นี่ไม่ใช่ QR token และใช้เงื่อนไขตารางเรียนเฉพาะทาง
\item \code{resolveRoomEntitlement} ใช้ legacy token เฉพาะที่ออกให้ user/room นั้น; group token ไม่ขยายสิทธิ์สู่สมาชิกใน resolver แม้ Gate validation ของ token กลุ่มมี branch แยก
\item การตอบ \code{valid: true} จาก server action/API Gate ไม่เท่ากับการบันทึก Session เสมอ; หน้า Gate จึงเรียก \code{checkInAction/checkOutAction} ต่อและตรวจผล. โค้ดไม่มีหลักฐานการควบคุมฮาร์ดแวร์ประตูแบบครบวงจร
\end{itemize}
\source{web/src/actions/access-tokens.ts; web/src/actions/groups.ts; web/src/app/(app)/dashboard/dashboard-client.tsx; web/src/app/gate/[roomCode]/page.tsx; web/src/app/api/gate/verify/route.ts}

\chapter{รหัสและ token สำคัญ}
\begin{longtable}{L{27mm}L{32mm}L{32mm}L{32mm}L{34mm}}
\toprule ชนิด & สร้างโดย/รูปแบบ & อายุ & เก็บ/ตรวจ & เพิกถอน/ข้อสังเกต \\ \midrule\endhead
Room code & Admin กำหนด; unique ใน Room & จนเปลี่ยน & DB, room lookup & ปิด room ได้ \\
Machine code & sequence ต่อห้อง, \code{-PCnn} & ตาม machine & unique ใน DB & deactivate machine \\
Batch activation & \code{randomInt}, 6 หลัก & 7 วัน & hash, consumedAt & revoke unused; ใช้ครั้งเดียว \\
Code เดี่ยว & \code{randomInt}, 6 หลัก & 10 นาที & hash, consumedAt & cancel unused \\
Machine API key & DB default/generated & จน deactivate & config เครื่อง + DB & deactivate/ถอน; ไม่ใส่ installer \\
Agent session token & \code{randomBytes(32)} base64url & ถึง entitlement end & SHA-256 hash ใน DB & logout/expiry/revoke \\
Maintenance code & \code{randomInt}, 6 หลัก & 5 นาที & hash, ผูก action/machine & ใช้ครั้งเดียว \\
Access token QR & \code{crypto.}\newline\code{randomUUID()} & validFrom--validTo & token ใน DB/QR & ไม่พบ consume เมื่อสแกน \\
Booking ID/batch ID & Prisma default \code{cuid()} & ตลอด record & primary key & สถานะ/ลบตามกฎ DB \\
\bottomrule
\end{longtable}
Hash activation/maintenance ใช้ helper \code{hashAgentOneTimeCode} ที่ผูก secret ของ server; session token เก็บ hash. API key ของ machine อยู่ใน config ที่ root/service อ่านได้และส่งใน header \code{x-agent-key}. Rate limit activation ใช้ durable throttle; Agent login/maintenance มีการจำกัดอัตราแยก. อย่านำ code ตัวอย่างในคู่มือไปใช้จริง
\source{web/src/lib/agent-activation.ts; web/src/lib/agent-deployment-batches.ts; web/src/lib/agent-session.ts; web/src/lib/agent-maintenance.ts; web/src/actions/access-tokens.ts; web/prisma/schema.prisma}

\chapter{Room Availability และ Entitlement}
Availability สำหรับการจองตรวจ room active, \code{bookingEnabled}, availability windows, closures, class schedule, booking เดิม, slot, capability, capacity, วันและเวลา. \code{AUTO} อนุมัติอัตโนมัติเมื่อผ่านกฎ; \code{MANUAL} ให้สถานะ pending รอ Admin. การจองแบบกลุ่มใช้ participant snapshot และ advisory locks กันการจองทับ/แข่งกัน; การเปลี่ยนสมาชิก Group ภายหลังไม่เปลี่ยนสิทธิ์ booking เดิม
\flowfour{Entitlement decision (ภาพอธิบาย)}{User + Room + Now}{สิทธิ์ 3 แหล่ง}{ตรวจช่วงเวลา}{Allow + validUntil}
\code{CLASS} ใช้ semester ที่ active, วันและช่วงเรียน โดยเริ่มก่อนชั้นเรียน 30 นาที. \code{BOOKING} ต้อง approved และ user อยู่ใน \code{RoomBookingParticipant}, เวลาอยู่ภายในช่วง slot. \code{LEGACY_TOKEN} ต้องออกให้ user และ room นี้และยังไม่หมดอายุ. หากหลายแหล่งทับกันให้สิทธิ์ถึงเวลาสิ้นสุดที่ไกลที่สุด. Agent login ใช้ resolver นี้; Face verify ใช้กับ STUDENT/GUEST. ตัวอย่างสมมติ: เรียน 09:00--12:00 ที่ห้อง LAB-A จะเริ่มสิทธิ์ CLASS 08:30 ถึงก่อน 12:00 ตามเวลา Asia/Bangkok
\source{web/src/lib/availability.ts; web/src/lib/room-booking.ts; web/src/lib/room-booking-rules.ts; web/src/lib/room-entitlement.ts; web/src/app/api/agent/auth/route.ts}

\chapter{Agent API และ Monitoring}
\begin{longtable}{L{50mm}L{19mm}L{43mm}L{47mm}}
\toprule Path & Method & ยืนยันตัวตน & หน้าที่/ข้อมูลสำคัญ \\ \midrule\endhead
\code{/api/agent/activate} & POST & code 6 หลัก + throttle & hostname/platform $\rightarrow$ machine code, room, API key \\
\code{/api/agent/auth} & POST & machine key + email/password & entitlement $\rightarrow$ sessionToken/expiresAt \\
\code{/api/agent/heartbeat} & POST & machine key; session optional & lastSeen/version; log activity เมื่อ session valid \\
\code{/api/agent/config} & GET & machine key & blacklist/room/platform \\
\code{/api/agent/session/refresh} & POST & machine key + session & ต่ออายุหาก entitlement ยังมี \\
\code{/api/agent/session/logout} & POST & machine key + session & revoke user session \\
\code{/api/agent/maintenance/}\newline\code{authorize} & POST & machine key + code & ตรวจ action; ถอนแล้ว deactivate \\
\code{/api/tracking/snapshot} & GET & Admin web session & สถานะเครื่อง/ผู้ใช้/กิจกรรมล่าสุด \\
\code{/api/tracking/machine-logs} & GET & Admin web session & log ล่าสุด 50 แถวตาม machineId \\
\code{/api/admin/agent/}\newline\code{installers/[platform]} & GET & Admin web session & EXE/PKG/DEB ที่ stage แล้ว \\
\bottomrule
\end{longtable}
Heartbeat ทุก 10 วินาทีอัปเดต \code{AgentMachine.lastSeen}; dashboard ถือ Online เมื่อ active และ lastSeen ใหม่กว่า 30 วินาที. Activity ล่าสุดแสดงเฉพาะเมื่อเครื่อง Online, มี Agent session ปัจจุบัน และ log เป็นผู้ใช้เดียวกันหลังเวลาออก session. Missing/invalid session token ส่ง 401; machine-only heartbeat ยังเป็น 200 และไม่สร้าง AgentLog. API ตัวอย่างไม่ควรฝัง key จริง
\section*{ข้อมูลและการตอบกลับที่ต้องรู้}
Activation รับ \code{activationCode}, \code{hostname}, \code{platform}; สำเร็จส่ง \code{machineCode}, \code{roomCode}, \code{apiKey}. รูปแบบผิดได้ 400, code ผิด/หมดอายุ/ใช้แล้ว/ถูกยกเลิกได้ 403 และ throttle ได้ 429 พร้อม \code{Retry-After}. Auth รับ \code{email/password} คู่ machine key; บัญชีไม่ approved หรือรหัสผิดได้ 401, ห้องไม่ active หรือไม่มี entitlement ได้ 403. Session token ต้องส่งใน \code{x-agent-session} เมื่อรายงาน activity. Maintenance รับ \code{code/action}; code ผิดได้ 403, machine key ผิดได้ 401 และ throttle ได้ 429. Snapshot และ machine logs ต้องมี Admin web session; installer route อาจตอบ 503 เมื่อยังไม่ตั้ง directory หรือ 404 เมื่อยังไม่ได้ stage file. ตัวอย่าง request ด้านล่างเป็นข้อมูลสมมติทั้งหมด
\begin{lstlisting}
POST /api/agent/heartbeat
X-Agent-Key: <machine-key-redacted>
Content-Type: application/json

{"agentVersion":"0.5.0"}
\end{lstlisting}
\source{web/src/app/api/agent/*/route.ts; web/src/app/api/tracking/snapshot/route.ts; web/src/lib/agent-health.ts; agent/src/api_client.py}

\chapter{ความปลอดภัย ความเป็นส่วนตัว และข้อจำกัด}
Backend แยก web session/role, machine key และ Agent session. Activation/maintenance code ถูก hash, ใช้ครั้งเดียว มีวันหมดอายุและ throttle. Machine config อยู่ในพื้นที่สิทธิ์จำกัด; UI ติดต่อ machine service ด้วย local token. HTTPS ใช้ที่ public endpoint. Admin ดู machine/user/activity ตามสิทธิ์; ข้อมูลชื่อหน้าต่างอาจมีข้อมูลส่วนตัว ควรกำหนดการแจ้งผู้ใช้ ระยะเก็บ และสิทธิ์เข้าถึงก่อนขยายใช้งานจริง

\begin{longtable}{L{47mm}L{107mm}}
\toprule สถานะ & ข้อเท็จจริง/งานที่ต้องตรวจ \\ \midrule\endhead
\status{IMPLEMENTED} & Agent machine heartbeat แยกจาก student login, QR gate path, Face AI match, room entitlement, batch codes, admin snapshot \\
\status{AUTOMATED TESTED} & รายงานก่อนหน้าระบุ Agent 19, Web 549, DB 166 และ build/CI ผ่าน; ไม่รันชุดนี้ซ้ำเพื่อทำเอกสาร \\
\status{VERIFIED ON REAL HARDWARE} & ผู้ใช้รายงาน macOS install/activate/Online/reboot UI และ machine identity เดิม รวมถึงถอนด้วย maintenance code และตรวจ receipt/launchctl \\
\status{NOT YET VERIFIED} & Student Agent login/activity/logout บนเครื่องจริง; Windows/Ubuntu install/reboot จริง; ประตูฮาร์ดแวร์; macOS title capture หลังให้สิทธิ์ \\
\status{PARTIAL IMPLEMENTATION} & เว็บไซต์ blacklist ยังไม่พิสูจน์การบังคับใช้ครบสาม OS; Group QR UI ปิด; Contact/Survey ไม่ควรนำเสนอเป็น workflow production \\
\status{NOT IMPLEMENTED} & Agent 0.5.0 ไม่บังคับ OS login ก่อนใช้ desktop; ไม่พบ liveness/anti-spoofing ใน Face flow \\
\bottomrule
\end{longtable}
คู่มือ Agent เก่าบางไฟล์ระบุ backend \code{/api/agent/activate} เป็น 404 ณ 11 พ.ย. 2026 ซึ่งเป็นหลักฐานประวัติ ไม่ใช่สถานะหลัง deployment ที่รายงานในวันถัดมา. การติดตั้ง/activate macOS ที่ผู้ใช้รายงานพิสูจน์ว่า endpoint ใช้งานได้ในช่วงทดลองนั้น แต่ไม่ได้ยืนยันทุก endpoint หรือทุก OS. การถอน macOS script ไม่ลบ log ใน \code{/var/log}; การออก maintenance code ถอนจะ deactivate machine ก่อนลบ local เสร็จ จึงต้องมีขั้นตรวจเมื่อเกิด partial failure
\source{agent/EXAM_CHEAT_SHEET.md; agent/macos/uninstall-agent.sh; web/src/app/api/agent/maintenance/authorize/route.ts; web/src/lib/agent-health.ts}

\appendix
\chapter{เช็กลิสต์ปฏิบัติงาน}
ทุกข้อเป็นเกณฑ์ที่ผู้ดูแลควรลงชื่อ/บันทึกผลในระบบงานของตนเอง ไม่ใช่หลักฐานว่าผ่านแล้ว คอลัมน์ผลให้ทำเครื่องหมายเมื่อมีหลักฐานจริง
\begin{longtable}{L{46mm}L{101mm}}
\toprule งาน & เกณฑ์ตรวจเสร็จ \\ \midrule\endhead
Backend ใหม่ & ตรวจ Node/pnpm/PM2/PostgreSQL, SHA release, .env โดยไม่เปิดเผยค่า, backup+restore test, migration status, build, smoke 3001, public health \\
Super Admin ครั้งแรก & บัญชี SUPER\_ADMIN ที่ได้รับอนุญาต login ได้; หน้า \code{/admin/system} แสดงเฉพาะ role นี้; บันทึกเจ้าของบัญชี \\
สร้างห้อง & roomCode ไม่ซ้ำ, active, capacity, capabilities, bookingEnabled, approvalPolicy, guestAccess/staffOnly, windows/closures ถูกต้อง \\
นำเข้าตาราง & semester/day/slot/room/user ตรงข้อมูลจริง; preview/รายการผิดพลาดตรวจแล้ว \\
Deployment Batch & ห้อง active; จำนวน 1--100; CSV ดาวน์โหลดครั้งเดียว; แจก code ตาม machine code; เก็บ CSV จำกัดสิทธิ์ \\
ติดตั้งหนึ่งเครื่อง & hash installer ตรง; OS/architecture ตรง; service/UI เริ่ม; activate; room/machine code ตรง; heartbeat Online \\
ติดตั้ง 30 เครื่อง & แบ่ง batch ตามห้อง; บันทึก asset/machine code; ตรวจไม่ใช้ code ซ้ำ; ตรวจ Online ครบและไม่มีเครื่องซ้ำ \\
Agent login & บัญชี approved มี entitlement ณ เวลานั้น; session ผูกเครื่อง/ห้อง/ผู้ใช้; logout แล้ว heartbeat ยังอยู่ \\
ทดสอบ OCR & ภาพสมมติที่อ่านได้; preview, warning, duplicate และ DB import ถูกต้อง \\
ทดสอบ Face & ลงทะเบียนด้วยผู้มีสิทธิ์; ตรวจ match/no match และ entitlement; ไม่อ้าง liveness/ประตู \\
ทดสอบ QR & UUID token ใน QR; gate สแกน; room/time validation; ตรวจผล check-in/out แยก \\
ถอน Windows & maintenance code ตรงเครื่อง/action; service/Run/ไฟล์หาย; backend inactive \\
ถอน macOS & receipt/LaunchDaemon/LaunchAgent หาย; backend inactive; ตรวจ log ที่ script ไม่ลบตามนโยบาย \\
ถอน Ubuntu & package/service/config ตาม postrm; backend inactive; legacy autostart ไม่เหลือ \\
อัปเดต Agent & hash/version ใหม่; config เดิมอยู่; service/UI เริ่ม; machine code เดิม; heartbeat กลับมา \\
สำรอง DB & \code{pg_dump -Fc}, SHA-256, \code{pg_restore --list}, restore แยกตามรอบทดสอบ \\
กู้ service & ระบุเหตุ/เวลา/ผู้รับผิดชอบ; rollback app ตาม runbook; ตรวจ PM2, Nginx, DB, HTTPS, Agent heartbeat \\
\bottomrule
\end{longtable}

\chapter{คำสั่งอ้างอิงแบบอ่านอย่างเดียว}
เลือกชุดคำสั่งตาม OS ของเครื่องที่จะตรวจ ไม่รัน installer/uninstaller บนเครื่องที่ไม่ใช่เป้าหมาย ไม่คัดลอกผลที่มี API key หรือข้อมูลส่วนบุคคลลง ticket สาธารณะ
\begin{lstlisting}[language=bash]
# macOS หลังติดตั้ง
sudo launchctl print system/com.elegantsolution.agent.machine
launchctl print gui/$(id -u)/com.elegantsolution.agent.session
pkgutil --pkg-info com.elegantsolution.agent

# Ubuntu หลังติดตั้ง
systemctl status elegant-agent.service --no-pager
systemctl is-enabled elegant-agent.service
journalctl -u elegant-agent.service -n 50 --no-pager

# Backend health แบบไม่ต้องมี credential
curl --fail --silent --show-error \
  https://elegant.fafromh0me.com/api/liveness
\end{lstlisting}
\begin{lstlisting}
# Windows PowerShell หลังติดตั้ง
Get-Service ElegantSolutionsMachine
sc.exe qc ElegantSolutionsMachine
Get-ItemProperty 'HKLM:\Software\Microsoft\Windows\CurrentVersion\Run' -Name ElegantAgentUI
\end{lstlisting}
ถ้า \code{launchctl print} หลังถอนรายงาน \code{Could not find service} และ \code{pkgutil} รายงาน \code{No receipt} นั่นสอดคล้องกับการถอนสำเร็จ; ตรวจ Admin Tracking แยกอีกครั้ง

\chapter{อภิธานศัพท์และสถานะ}
\begin{longtable}{L{41mm}L{106mm}}
\toprule คำ & ความหมายในระบบนี้ \\ \midrule\endhead
Web session & ตัวตนผู้ใช้เว็บจาก NextAuth JWT; แยกจาก Agent session \\
Machine identity & machine code/API key หลัง activation; ใช้ machine heartbeat \\
Agent session & token ที่ผูกผู้ใช้-เครื่อง-ห้องและหมดอายุตาม entitlement \\
Entitlement & สิทธิ์ใช้ห้อง ณ เวลาใช้งานจาก CLASS, BOOKING หรือ LEGACY\_TOKEN \\
Heartbeat & ข้อความจาก service เพื่ออัปเดต lastSeen; ไม่ยืนยันผู้ใช้ \\
Online & เครื่อง active และ lastSeen ภายใน 30 วินาที \\
Deployment Batch & ชุด code ที่จองให้ห้องและจำนวนเครื่อง; code ใช้ครั้งเดียว \\
Maintenance code & code หกหลักผูกเครื่อง/action มีอายุ 5 นาที \\
Gate & หน้าเว็บ/endpoint ตรวจ QR หรือใบหน้า; ไม่ใช่หลักฐานฮาร์ดแวร์เปิด \\
Preview & รายการ OCR ก่อนบันทึกที่ผู้ใช้ต้องตรวจ \\
\bottomrule
\end{longtable}

\chapter{บัญชีประเด็นเอกสารและแหล่งอ้างอิง}
\section*{ข้อแตกต่างและสิ่งที่ควรตรวจต่อ}
\begin{enumerate}
\step{\code{agent/*INSTALLATION.md} และ \code{agent/EXAM_CHEAT_SHEET.md} ยังกล่าวถึง production endpoint 404 เมื่อ 11 พ.ย. 2026; สถานะนั้นเก่ากว่าผลการติดตั้ง/activate macOS ที่ผู้ใช้รายงาน}
\step{Agent 0.5.0 ไม่บังคับ OS login ก่อน desktop ใช้งาน; หากต้องการควบคุมเครื่องแล็บจริงต้องออกแบบ OS-level policy และทดสอบบนเครื่องที่จัดการได้}
\step{\code{web/prisma/seed.ts} สร้างบัญชีตัวอย่าง ADMIN พร้อม credential ที่ฝังใน source แต่ไม่มี SUPER\_ADMIN; ต้องวางขั้น provision ครั้งแรกที่ปลอดภัยก่อนใช้ระบบจริง และห้ามรัน seed นี้บน production}
\step{QR access token legacy และ room entitlement ใหม่มีเงื่อนไขไม่เหมือนกัน; ทบทวน policy ก่อนขยายการใช้งาน Gate จริง}
\step{Room Finder ซ่อนห้อง \code{staffOnly} จาก Guest แต่ \code{generateAccessTokenAction} ตรวจ \code{guestAccess} โดยไม่ตรวจ \code{staffOnly} โดยตรง; ทบทวนการบังคับสิทธิ์ Gate/token ก่อนใช้งานจริง}
\step{Face AI เลือกหน้าที่ detection score สูงสุดเมื่อภาพมีหลายหน้า; ยังไม่มี liveness/anti-spoofing ใน flow ที่ตรวจ}
\step{macOS uninstaller ไม่ลบ log สองไฟล์ใน \code{/var/log}; backend deactivate machine ในขั้น authorize ก่อน local cleanup จบ}
\step{Windows/Ubuntu ยังต้องติดตั้งและทดสอบ reboot/activity จริง; macOS foreground activity และ student login ก็ยังไม่มีผลทดสอบที่ส่งมา}
\end{enumerate}
\section*{แผนที่หลักฐานใน repository}
\begin{itemize}
\item โครงสร้างและ role: \code{web/prisma/schema.prisma}, \code{web/src/lib/permissions.ts}, \code{web/src/lib/auth.ts}
\item ห้องและการจอง: \code{web/src/lib/availability.ts}, \code{web/src/lib/room-booking.ts}, \code{web/src/lib/room-entitlement.ts}
\item OCR/Face/Gate: \code{web/src/lib/ocr-timetable.ts}, \code{web/src/app/api/face/*}, \code{web/src/actions/access-tokens.ts}
\item Agent: \code{agent/src/machine_service.py}, \code{agent/src/session_ui.py}, \code{agent/build-macos.sh}, \code{agent/build-linux.sh}, \code{agent/build.bat}
\item Deployment: \code{docs/PRODUCTION_PM2_DEPLOYMENT.md}; Agent API: \code{web/src/app/api/agent/*}
\end{itemize}
\note{ภาพทั้งหมดในคู่มือเป็นแผนภาพต้นฉบับจากโค้ด ไม่ใช่ภาพหน้าจอ production; ตัวอย่างรหัส/ข้อมูลใช้ข้อความสมมติ ไม่มี backup, secret หรือข้อมูลใบหน้าจริงอยู่ในเอกสาร}

\end{document}
